% 第三章：曲线积分

\section{曲线积分}

数二的积分对象是区间（定积分）和区域（二重积分），积分的"舞台"是直线段或平面区域。数一引入了曲线积分，将积分的舞台推广到了曲线——一条弯曲的一维对象。曲线积分分为两类：第一类曲线积分（对弧长的积分）和第二类曲线积分（对坐标的积分），它们分别对应着不同的物理原型和几何意义。

\subsection{第一类曲线积分：弧上的"称重"}

\subsubsection{定义与物理原型}

设 $L$ 是一条光滑曲线弧，$f(x,y)$ 在 $L$ 上有界。将 $L$ 分成 $n$ 个小弧段 $\Delta s_i$，在第 $i$ 段上任取一点 $(\xi_i,\eta_i)$，作和式 $\sum_{i=1}^n f(\xi_i,\eta_i)\Delta s_i$。当最大弧长 $\lambda\to 0$ 时，若极限存在，则称为 $f$ 在 $L$ 上对弧长的曲线积分：
\[
\int_L f(x,y)\,\mathrm{d}s = \lim_{\lambda\to 0}\sum_{i=1}^n f(\xi_i,\eta_i)\Delta s_i
\]

\begin{essencebox}
第一类曲线积分的本质是\textbf{沿曲线的加权累积}。如果 $f(x,y)$ 表示曲线 $L$ 上点 $(x,y)$ 处的线密度，则 $\int_L f\,\mathrm{d}s$ 就是曲线的总质量。微元 $\mathrm{d}s$ 是弧长元素，代表曲线上一个"无穷小"的弧段长度。当 $f\equiv 1$ 时，积分结果就是曲线的弧长。

关键点：第一类曲线积分与曲线的\textbf{方向无关}。无论从 $A$ 到 $B$ 还是从 $B$ 到 $A$ 积分，结果相同，因为弧长 $\mathrm{d}s$ 始终为正。
\end{essencebox}

\subsubsection{计算方法}

若曲线 $L$ 的参数方程为 $x=x(t), y=y(t)$，$t\in[\alpha,\beta]$，则：
\[
\int_L f(x,y)\,\mathrm{d}s = \int_\alpha^\beta f(x(t),y(t))\sqrt{x'(t)^2+y'(t)^2}\,\mathrm{d}t
\]

\begin{insightbox}
\textbf{推理步骤}：为什么出现 $\sqrt{x'^2+y'^2}$？
\begin{enumerate}[leftmargin=2em]
    \item 弧长微元 $\mathrm{d}s$ 是曲线上无穷小弧段的长度
    \item 由勾股定理，$\mathrm{d}s = \sqrt{\mathrm{d}x^2+\mathrm{d}y^2}$
    \item 参数化后，$\mathrm{d}x = x'(t)\,\mathrm{d}t$，$\mathrm{d}y = y'(t)\,\mathrm{d}t$
    \item 因此 $\mathrm{d}s = \sqrt{x'(t)^2+y'(t)^2}\,\mathrm{d}t$
    \item 代入积分即得公式
\end{enumerate}
这个因子 $\sqrt{x'^2+y'^2}$ 的本质是参数 $t$ 的"速度"——参数变化一个单位，曲线上的点移动了多远。
\end{insightbox}

\subsubsection{空间曲线的情形}

对于空间曲线 $\Gamma: x=x(t), y=y(t), z=z(t)$，$t\in[\alpha,\beta]$，则：
\[
\int_\Gamma f(x,y,z)\,\mathrm{d}s = \int_\alpha^\beta f(x(t),y(t),z(t))\sqrt{x'(t)^2+y'(t)^2+z'(t)^2}\,\mathrm{d}t
\]

空间曲线的情形与平面曲线完全类似，只是多了一个 $z$ 分量。弧长微元变为 $\mathrm{d}s = \sqrt{\mathrm{d}x^2+\mathrm{d}y^2+\mathrm{d}z^2}$，这正是三维空间中两点间距离公式的微分形式。

\subsection{第二类曲线积分：力沿路径做功}

\subsubsection{定义与物理原型}

第二类曲线积分的定义为：
\[
\int_L P\,\mathrm{d}x + Q\,\mathrm{d}y \quad \text{或} \quad \int_L \vec{F}\cdot\mathrm{d}\vec{r}
\]

其中 $\vec{F} = (P, Q)$ 是向量场，$\mathrm{d}\vec{r} = (\mathrm{d}x, \mathrm{d}y)$ 是曲线的有向弧长元素。

\begin{essencebox}
第二类曲线积分的物理原型是\textbf{变力沿曲线做功}。设质点在力场 $\vec{F}$ 的作用下沿曲线 $L$ 从 $A$ 移动到 $B$，则力场做的功为 $W = \int_L \vec{F}\cdot\mathrm{d}\vec{r}$。微元 $\vec{F}\cdot\mathrm{d}\vec{r}$ 是力在位移方向上的分量乘以位移，即微元功。

关键点：第二类曲线积分与曲线的\textbf{方向有关}。从 $A$ 到 $B$ 和从 $B$ 到 $A$ 的积分结果相差一个负号，因为做功与力的方向和位移方向都有关。
\end{essencebox}

\subsubsection{计算方法}

若曲线 $L$ 的参数方程为 $x=x(t), y=y(t)$，$t$ 从 $\alpha$ 到 $\beta$，则：
\[
\int_L P\,\mathrm{d}x + Q\,\mathrm{d}y = \int_\alpha^\beta \left[P(x(t),y(t))x'(t) + Q(x(t),y(t))y'(t)\right]\mathrm{d}t
\]

\begin{insightbox}
\textbf{推理步骤}：为什么第二类曲线积分的公式与第一类不同？
\begin{enumerate}[leftmargin=2em]
    \item 第一类积分的微元是标量 $\mathrm{d}s$（弧长），始终为正
    \item 第二类积分的微元是向量 $\mathrm{d}\vec{r} = (\mathrm{d}x, \mathrm{d}y)$，方向由曲线方向决定
    \item 参数化后，$\mathrm{d}x = x'(t)\,\mathrm{d}t$，$\mathrm{d}y = y'(t)\,\mathrm{d}t$
    \item 当 $t$ 从 $\alpha$ 增加到 $\beta$ 时，参数增加的方向就是曲线的正方向
    \item 因此 $\mathrm{d}t$ 的符号自动处理了方向问题
\end{enumerate}
\end{insightbox}

\subsubsection{两类曲线积分的关系}

两类曲线积分通过单位切向量 $\vec{\tau}$ 联系起来。设 $\vec{\tau} = (\cos\alpha, \cos\beta)$ 是曲线 $L$ 在点 $(x,y)$ 处的单位切向量，则：
\[
\int_L P\,\mathrm{d}x + Q\,\mathrm{d}y = \int_L (P\cos\alpha + Q\cos\beta)\,\mathrm{d}s
\]

即第二类曲线积分等于被积函数与切向量的点积的第一类曲线积分。这个关系揭示了两类积分的本质区别：第一类积分是"无方向的加权累积"，第二类积分是"有方向的加权累积"。

\subsection{格林公式：曲线积分与二重积分的桥梁}

\subsubsection{格林公式}

设 $D$ 是由分段光滑闭曲线 $L$ 围成的闭区域，$P(x,y)$ 和 $Q(x,y)$ 在 $D$ 上有连续一阶偏导数，则：
\[
\oint_L P\,\mathrm{d}x + Q\,\mathrm{d}y = \iint_D \left(\frac{\partial Q}{\partial x} - \frac{\partial P}{\partial y}\right)\mathrm{d}\sigma
\]

其中 $L$ 取正方向（逆时针方向）。

\begin{essencebox}
格林公式的本质是\textbf{微积分基本定理在二维的推广}。牛顿-莱布尼茨公式 $\int_a^b f'(x)\,\mathrm{d}x = f(b)-f(a)$ 将区间端点的值与区间内的积分联系起来；格林公式将边界曲线上的积分与区域内的二重积分联系起来。这种"边界积分 = 内部积分"的模式是微积分基本定理的核心，在三维中对应高斯公式和斯托克斯公式。
\end{essencebox}

\begin{insightbox}
\textbf{推理步骤}：格林公式的证明思路
\begin{enumerate}[leftmargin=2em]
    \item 只需证明 $\oint_L P\,\mathrm{d}x = -\iint_D \frac{\partial P}{\partial y}\,\mathrm{d}\sigma$ 和 $\oint_L Q\,\mathrm{d}y = \iint_D \frac{\partial Q}{\partial x}\,\mathrm{d}\sigma$
    \item 对第一个等式：右端 $= -\int_a^b\int_{y_1(x)}^{y_2(x)}\frac{\partial P}{\partial y}\,\mathrm{d}y\,\mathrm{d}x = -\int_a^b[P(x,y_2)-P(x,y_1)]\,\mathrm{d}x$
    \item 左端：将 $L$ 分为上边界 $L_2$（从右到左）和下边界 $L_1$（从左到右）
    \item 在 $L_1$ 上 $y=y_1(x)$，$\mathrm{d}y=0$；在 $L_2$ 上方向相反
    \item 化简后两端相等，类似证明第二个等式
\end{enumerate}
\end{insightbox}

\subsubsection{格林公式的应用}

\textbf{1. 简化曲线积分计算}

当直接计算曲线积分困难，但 $\frac{\partial Q}{\partial x} - \frac{\partial P}{\partial y}$ 较简单时，可用格林公式将曲线积分转化为二重积分。

\textbf{2. 计算平面区域面积}

令 $P=-y$，$Q=x$，则 $\frac{\partial Q}{\partial x}-\frac{\partial P}{\partial y}=2$，由格林公式：
\[
A = \frac{1}{2}\oint_L x\,\mathrm{d}y - y\,\mathrm{d}x
\]

\textbf{3. 判定曲线积分与路径无关}

若 $\frac{\partial Q}{\partial x} = \frac{\partial P}{\partial y}$ 在单连通区域 $D$ 内处处成立，则 $\int_L P\,\mathrm{d}x + Q\,\mathrm{d}y$ 在 $D$ 内与路径无关。这是格林公式的重要推论。

\subsection{如何促进其他部分的理解}

\begin{connectionbox}
\textbf{1. 从牛顿-莱布尼茨到格林：微积分基本定理的统一}

牛顿-莱布尼茨公式将一维区间上的积分与端点的值联系起来；格林公式将二维区域上的积分与边界上的曲线积分联系起来。这种"内部积分 = 边界积分"的模式在后续的高斯公式（三维）和斯托克斯公式中将继续出现，构成了微积分基本定理的统一框架。

\textbf{2. 为曲面积分提供类比}

第一类曲线积分（对弧长）和第二类曲线积分（对坐标）的区分，在曲面积分中有完全平行的对应：第一类曲面积分（对面积）和第二类曲面积分（对坐标）。理解了曲线积分的两类区分，就能自然地理解曲面积分的两类区分。

\textbf{3. 为斯托克斯公式铺设道路}

格林公式是斯托克斯公式在平面上的特例。斯托克斯公式将空间曲线积分与曲面积分联系起来，而格林公式将平面曲线积分与二重积分联系起来。理解格林公式的本质，是理解斯托克斯公式的必要前提。

\textbf{4. 深化对"保守场"的理解}

格林公式揭示了"曲线积分与路径无关"的充要条件是 $\frac{\partial Q}{\partial x}=\frac{\partial P}{\partial y}$，即向量场"无旋"。这为理解保守场（势场）提供了二维的直观基础，在三维中对应旋度为零的条件。
\end{connectionbox}
